% =====================================================================
%  backmatter/lampiran-c.tex
%  Lampiran C: Kode Sumber Utama dan Repositori Data Penelitian.
% =====================================================================
\lampiran[lamp:kode]{Kode Sumber Utama dan Repositori Data Penelitian}

\panduan{Ketentuan Repositori dan Kode Sumber:
Sesuai standar riset program studi, mahasiswa WAJIB mengunggah:
(1) Kode sumber lengkap ke repositori publik GitHub, dilengkapi README.md dan dokumentasi instalasi;
(2) Dataset mentah, hasil pengolahan, dan metadata riset ke repositori data ilmiah resmi (Mendeley Data atau Zenodo) sehingga memiliki persistent identifier (DOI).
Pada naskah cetak skripsi, mahasiswa hanya mencantumkan tautan repositori serta potongan kode algoritma inti, bukan seluruh baris kode aplikasi.}

\noindent Informasi repositori resmi untuk penelitian ini adalah sebagai berikut:
\begin{enumerate}[noitemsep,topsep=3pt,leftmargin=1.5em]
  \item \textbf{Repositori Kode Sumber Utama (GitHub):}\\
        \url{\tautanrepogithub}
  \item \textbf{Repositori Dataset dan Artefak Penelitian (Mendeley Data):}\\
        \url{\tautanrepomendeley}
\end{enumerate}

\vspace{0.8\baselineskip}
\noindent Potongan kode program utama untuk perhitungan metrik evaluasi kegunaan sistem disajikan pada Kode~\ref{kode:sus}, yang mengimplementasikan Algoritma~\ref{alg:sus}.

\begin{lstlisting}[language=Python,caption={Implementasi Algoritma Perhitungan Skor SUS pada Backend Aplikasi},label={kode:sus}]
def hitung_skor_sus(jawaban_responden):
    """
    Menghitung skor System Usability Scale (SUS) skala 0-100
    berdasarkan 10 butir pertanyaan dengan skala Likert 1-5.
    """
    if len(jawaban_responden) != 10:
        raise ValueError("Instrumen SUS mewajibkan tepat 10 butir respons.")
    
    total_skor = 0
    for nomor_butir, nilai in enumerate(jawaban_responden, start=1):
        if nomor_butir % 2 == 1:
            # Butir bernomor ganjil (pernyataan positif): nilai - 1
            total_skor += (nilai - 1)
        else:
            # Butir bernomor genap (pernyataan negatif): 5 - nilai
            total_skor += (5 - nilai)
            
    # Nilai akhir dikalikan 2.5 untuk konversi rentang 0-100
    return float(total_skor * 2.5)
\end{lstlisting}
